\section{Introduction}
\label{sec:introduction}

Titanium aluminides combine low density with high specific strength and creep resistance at elevated temperatures, making them attractive for lightweight components in aircraft and automotive engines.\cite{Clemens2013}
The binary Ti--Al system provides the chemical basis for these alloys and contains several ordered compounds, including $\gamma$-TiAl and $\alpha_2$-Ti$_3$Al.
Their processing and performance involve both crystalline phases and the liquid state.
In particular, simulations of melting and solidification require an accurate description of phase stability, while the density, mixing thermodynamics, and viscosity of the melt provide complementary information about its behavior during processing.
Developing an atomistic description that remains reliable across compositions and phases is therefore an important problem for Ti--Al alloys.

Considerable effort has already been devoted to constructing interatomic potentials for this system.
Established approaches include the embedded-atom potentials of Zope and Mishin\cite{Zope2003} and the second-nearest-neighbor modified embedded-atom potential of Kim \textit{et al.}\cite{Kim2016}
Machine-learning interatomic potentials (MLIPs) have subsequently expanded the flexibility of the representation and the range of reference configurations that can be included in training.
Examples include the polynomial-invariant potentials of Seko for binary Ti--Al,\cite{Seko2020} the neuroevolution potential developed by Zhao \textit{et al.} for Ti--Al--Nb,\cite{Zhao2024} and the Ti--Al neuroevolution potential of Shen \textit{et al.}\cite{Shen2025}
Recent work has also examined the transferability of broadly trained potentials to Al--Ti melts, including their composition-dependent thermophysical and transport properties.\cite{Song2024,Kato2026}
These studies establish a substantial basis for atomistic modeling of Ti--Al and make the accuracy of the underlying reference description an increasingly relevant question.

The three system-specific MLIP studies cited above use density-functional theory (DFT) reference calculations with the Perdew--Burke--Ernzerhof (PBE) exchange-correlation functional.\cite{Seko2020,Zhao2024,Shen2025,Perdew1996}
Agreement with such reference data measures how accurately a potential reproduces the selected DFT approximation; agreement with experiment also depends on the physical accuracy of that approximation.
The melting temperature of elemental Ti provides a particularly clear example.
Zhao \textit{et al.} report a value of $1671\pm5$~K, compared with their experimental reference of 1941~K,\cite{Zhao2024} whereas Shen \textit{et al.} obtain 1611~K against 1943~K.\cite{Shen2025}
Both models reproduce a range of crystalline properties, yet their Ti melting temperatures are lower than experiment by approximately 14--17\%.
These discrepancies motivate an explicit examination of exchange-correlation effects.
Differences in training coverage, fitting errors, and simulation protocols must also be considered when comparing the predictions of distinct MLIPs.

Melting and liquid transport are demanding tests because they probe aspects of the potential-energy surface that are not fully characterized by equilibrium crystal properties.
The melting temperature depends on the free-energy balance between the solid and liquid, and viscosity depends on the relaxation of microscopic stress fluctuations.
Consequently, small regression errors for energies and forces do not by themselves establish quantitative accuracy for either observable.
First-principles calculations accelerated by machine-learned force fields have shown that predicted melting temperatures depend appreciably on the exchange-correlation approximation.\cite{Jinnouchi2019}
For alloy liquids, this dependence can additionally vary with composition, making a comparison over the binary composition range more informative than a test of the pure elements alone.

Our recent study of Al--Cu melts demonstrated the practical importance of this issue.\cite{Khazieva2024}
PBE-based potentials described a broad set of properties but failed to reproduce the measured composition dependence of viscosity in the Cu-rich region.
Fine-tuning on reference data calculated with PBEsol or \rscan{} substantially improved the description of thermodynamic and transport properties; the \rscan{} potential reproduced the experimentally observed viscosity maximum.
That work also showed that transfer learning can reduce the amount of reference data needed to construct potentials at a different level of electronic-structure theory.
These results motivate a corresponding assessment for Ti--Al, where the balance between crystalline and liquid properties, and its variation with composition, remains to be established for alternative functionals.

PBEsol and \rscan{} provide two distinct alternatives to PBE for this comparison.
PBEsol modifies the generalized-gradient approximation to improve the description of densely packed solids and their surfaces,\cite{Perdew2008} whereas \rscan{} is a regularized, constraint-restored meta-generalized-gradient approximation.\cite{Furness2020}
Their different constructions offer an opportunity to examine how changes in the reference electronic-structure approximation affect crystalline and liquid observables within a common MLIP framework.
Improved agreement for one property need not imply an equivalent improvement for the others, so the assessment must include several independent physical benchmarks.

Foundation interatomic models make such comparisons more accessible.
Pretraining on broad chemical and structural datasets provides representations that can be adapted to a particular system using a comparatively compact reference dataset.\cite{Batatia2025}
Transfer learning can then be used to change both the material domain and the electronic-structure level represented by the potential.\cite{Khazieva2024}
The recently introduced DPA4 architecture provides an equivariant, energy-conserving framework for this purpose.\cite{Li2026DPA4}
Its forces are obtained as derivatives of the predicted energy, and pretrained DPA4 models are available in several sizes.
Here we use DPA4-Neo as the common starting model for fine-tuning.
Using the same pretrained backbone and configuration-selection strategy provides a consistent basis for examining the effect of changing the DFT reference functional.

In this work, we develop Ti--Al MLIPs using PBE, PBEsol, and \rscan{} reference calculations and compare their predictions for crystalline and liquid states.
The reference configurations are selected from a common pool generated by foundation-model molecular dynamics and structural relaxation, and are used to fine-tune the DPA4-Neo model.
We assess lattice dynamics, temperature-dependent density, the melting temperature of Ti, and the composition dependence of liquid density and mixing enthalpy.
Liquid viscosity provides an additional comparison between the PBE- and \rscan{}-based models.
The combined benchmarks allow us to determine which improvements persist across observables and which depend on composition or phase.
The resulting assessment addresses the choice of electronic-structure reference data for Ti--Al MLIPs and the use of foundation-model fine-tuning as a practical route to that comparison.
